\documentclass[11pt]{article}
\usepackage[margin=1in]{geometry}
\usepackage{enumitem}
% =============================================================================
% Preambles/header.tex — shared LaTeX/Beamer preamble for this project
%
% Use from a Beamer lecture by prepending:
%     \input{header}
% and compiling with TEXINPUTS=../Preambles:$TEXINPUTS (the /compile-latex
% skill does this automatically).
%
% PALETTE CONTRACT. Color names defined here MUST match the names in
% Quarto/theme-template.scss so Beamer and Quarto renderings use the same
% palette. The scripts/check-palette-sync.sh script enforces this.
%
% To customize: edit the HEX values below AND the matching $variable in
% Quarto/theme-template.scss, then run scripts/check-palette-sync.sh.
% =============================================================================

% -----------------------------------------------------------------------------
% Engine detection + encoding
%
% Our /compile-latex skill uses XeLaTeX, where inputenc is unnecessary and
% can produce encoding warnings. Only load inputenc under pdfTeX.
% -----------------------------------------------------------------------------
\usepackage{iftex}
\ifPDFTeX
  \usepackage[utf8]{inputenc}
\fi

\usepackage{amsmath, amssymb, amsthm}
\usepackage{booktabs}
\usepackage{graphicx}

% -----------------------------------------------------------------------------
% PALETTE — mirrors Quarto/theme-template.scss variable names.
% Load xcolor and define colors BEFORE anything that references them
% (notably hyperref, hypersetup, and the Beamer color assignments).
% -----------------------------------------------------------------------------
\usepackage{xcolor}

\definecolor{primary-blue}{HTML}{012169}      % headings, accents
\definecolor{primary-gold}{HTML}{B9975B}      % emphasis, borders
\definecolor{highlight-yellow}{HTML}{F2A900}  % markers, alerts
\definecolor{light-bg}{HTML}{E8EDF5}          % title / section backgrounds
\definecolor{jet}{HTML}{1A1A1A}               % body text

% Semantic colors (used in diagrams and annotations)
\definecolor{positive}{HTML}{15803D}          % good / observed / identified
\definecolor{negative}{HTML}{B91C1C}          % bad / problematic / bias
\definecolor{neutral}{HTML}{525252}           % reference / context

% Highlight accents (match .hi-* classes in the SCSS)
\definecolor{hi-slate}{HTML}{314F4F}
\definecolor{hi-green}{HTML}{2E8B57}
\definecolor{hi-red}{HTML}{C41E3A}

% Semantic aliases so skill templates and snippets can write
% \textcolor{accent}{...} without hard-coding "primary-blue".
\colorlet{accent}{primary-blue}
\colorlet{accent2}{primary-gold}

% -----------------------------------------------------------------------------
% Hyperref — loaded AFTER xcolor + palette so link colors resolve.
% -----------------------------------------------------------------------------
\usepackage{hyperref}
\hypersetup{colorlinks=true, linkcolor=primary-blue, urlcolor=primary-blue,
            citecolor=primary-blue, pdfborder={0 0 0}}

% -----------------------------------------------------------------------------
% Course code — set once per deck (after \input{header}, before \begin{document})
% via \coursecode{CS401}. Shown in the footer on every slide so a deck's course
% is identifiable without opening the title page. Empty by default.
% -----------------------------------------------------------------------------
\makeatletter
\newcommand{\coursecode}[1]{\gdef\@coursecodetext{#1}}
\gdef\@coursecodetext{}
\makeatother

% -----------------------------------------------------------------------------
% Beamer theme assignments (only apply under Beamer).
%
% We detect Beamer via \@ifclassloaded{beamer}, which is the documented,
% non-internal way. This must be wrapped in \makeatletter/\makeatother
% because \@ifclassloaded contains an @-catcode character.
% -----------------------------------------------------------------------------
\makeatletter
\@ifclassloaded{beamer}{%
  \usefonttheme{professionalfonts}
  \setbeamercolor{structure}{fg=primary-blue}
  \setbeamercolor{frametitle}{fg=primary-blue}
  \setbeamercolor{title}{fg=primary-blue}
  \setbeamercolor{subtitle}{fg=primary-gold}
  \setbeamercolor{author}{fg=primary-blue}
  \setbeamercolor{institute}{fg=neutral}
  \setbeamercolor{date}{fg=primary-gold}
  \setbeamercolor{itemize item}{fg=primary-blue}
  \setbeamercolor{itemize subitem}{fg=primary-gold}
  \setbeamercolor{alerted text}{fg=primary-gold}
  \setbeamercolor{block title}{fg=white, bg=primary-blue}
  \setbeamercolor{block body}{bg=light-bg}
  % Semantic block variants: block=definition/concept (blue), exampleblock=
  % worked example (green/positive), alertblock=key takeaway or caution (gold).
  % Keep to <=2 colored boxes per slide across ALL three variants (INV-7).
  \setbeamercolor{block title example}{fg=white, bg=positive}
  \setbeamercolor{block body example}{bg=light-bg}
  \setbeamercolor{block title alerted}{fg=white, bg=primary-gold}
  \setbeamercolor{block body alerted}{bg=light-bg}

  % Minimal footer: course code (left) + page X / N (right), no navigation clutter
  \setbeamertemplate{navigation symbols}{}
  \setbeamertemplate{footline}{%
    \color{neutral}\footnotesize\hspace{0.7em}\@coursecodetext\hfill
    \insertframenumber/\inserttotalframenumber\hspace{0.7em}\vspace{0.3em}}
}{}
\makeatother

% -----------------------------------------------------------------------------
% TikZ setup — libraries the snippet gallery uses
% -----------------------------------------------------------------------------
\usepackage{tikz}
\usetikzlibrary{arrows.meta, positioning, calc, decorations.pathreplacing,
                fit, shapes.geometric, backgrounds}

% Shared TikZ styles — snippets and hand-written diagrams can pick these up
% via `\begin{tikzpicture}[every node/.style={...}]` or by naming specific
% styles (e.g., `\node[dag-node] (X) at (0,0) {X};`).
\tikzset{
  % Node styles for causal DAGs and similar
  dag-node/.style = {
    draw=primary-blue, fill=light-bg, rounded corners=2pt,
    minimum width=1.2cm, minimum height=0.9cm, align=center, font=\small
  },
  decision-node/.style = {
    draw=primary-gold, fill=white, diamond, aspect=1.8,
    minimum width=1.8cm, minimum height=1.2cm, align=center, font=\small,
    inner sep=1pt
  },
  % Edge styles — observed vs counterfactual
  observed-edge/.style = {-{Latex[length=2.5mm]}, thick, draw=jet},
  counterfactual-edge/.style = {-{Latex[length=2.5mm]}, thick, draw=neutral, dashed},
  confound-edge/.style = {-{Latex[length=2.5mm]}, thick, draw=negative, dashed},
  % Data-point styles for plots
  observed-dot/.style = {fill=primary-blue, draw=primary-blue, circle, inner sep=1.2pt},
  counterfactual-dot/.style = {fill=white, draw=neutral, circle, inner sep=1.2pt}
}

% -----------------------------------------------------------------------------
% Convenience macros
% -----------------------------------------------------------------------------
\newcommand{\muted}[1]{\textcolor{neutral}{#1}}
\newcommand{\key}[1]{\textcolor{primary-gold}{\textbf{#1}}}
\newcommand{\good}[1]{\textcolor{positive}{#1}}
\newcommand{\bad}[1]{\textcolor{negative}{#1}}

% Standout frame macro: full-bleed dark background for section transitions.
% Beamer-only (uses \begin{frame}/\setbeamercolor/\usebeamerfont) -- do not
% call from an article-class file (e.g. Notes/<CODE>/*-notes.tex). \muted,
% \key, \good, \bad, and \coursecode above are plain xcolor/gdef and are
% safe to use from either class; verified when Notes/article-class support
% was added.
% Expand: \transitionslide{Your transition title}
\newcommand{\transitionslide}[1]{%
  \begingroup
  \setbeamercolor{background canvas}{bg=primary-blue}
  \begin{frame}[plain]
    \centering
    \vfill
    {\usebeamerfont{title}\color{white}\Large #1\par}
    \vfill
  \end{frame}
  \endgroup
}

% Section divider frame (Beamer-only), an alternative to \transitionslide with a
% darker two-line style: near-black (jet) background, a small white label line
% above a larger gold title, and a thin gold rule along the bottom edge.
% Expand: \sectiondivider{Section 2}{Pointers}
\newcommand{\sectiondivider}[2]{%
  \begingroup
  \setbeamercolor{background canvas}{bg=jet}
  \begin{frame}[plain]
    \vspace*{3.0cm}
    {\color{white}\ttfamily\large #1\par}
    \vspace{0.5cm}
    {\color{primary-gold}\LARGE #2\par}
    \begin{tikzpicture}[remember picture, overlay]
      \draw[primary-gold, line width=2.5pt]
        ($(current page.south west)+(0,0.1)$) -- ($(current page.south east)+(0,0.1)$);
    \end{tikzpicture}
  \end{frame}
  \endgroup
}

% -----------------------------------------------------------------------------
% Channel watermark — "Concepts That Click" (YouTube).
%
% Article-class files (CompetitiveExam Q/A sets, Notes, Assignments) get a
% light diagonal draft watermark on every page plus a centered footer naming
% the channel. Beamer decks get a subtle TikZ background watermark instead
% (the footline already carries the course code). Centralized here so every
% MCQ PDF is branded without per-file edits.
% -----------------------------------------------------------------------------
\makeatletter
\@ifclassloaded{article}{%
  \usepackage{draftwatermark}
  \SetWatermarkText{Concepts That Click}
  \SetWatermarkScale{1}
  \SetWatermarkFontSize{2cm}
  \SetWatermarkLightness{0.86}
  \SetWatermarkAngle{45}
  \usepackage{fancyhdr}
  \pagestyle{fancy}
  \fancyhf{}
  \fancyfoot[C]{\footnotesize\textcolor{neutral}{Concepts That Click~|~YouTube: \href{https://www.youtube.com/@concepts.thatclick}{@concepts.thatclick}}}
  \renewcommand{\headrulewidth}{0pt}
  \renewcommand{\footrulewidth}{0pt}
  \fancypagestyle{plain}{%
    \fancyhf{}%
    \fancyfoot[C]{\footnotesize\textcolor{neutral}{Concepts That Click~|~YouTube: \href{https://www.youtube.com/@concepts.thatclick}{@concepts.thatclick}}}%
    \renewcommand{\headrulewidth}{0pt}%
    \renewcommand{\footrulewidth}{0pt}%
  }
}{}
\@ifclassloaded{beamer}{%
  \setbeamertemplate{background}{%
    \begin{tikzpicture}[remember picture, overlay]
      \node[rotate=45, opacity=0.055, align=center]
        at (current page.center)
        {\Huge\textcolor{neutral}{Concepts That Click}};
    \end{tikzpicture}%
  }
}{}
\makeatother 
\coursecode{JKSSB}
\renewcommand{\thesection}{55.\arabic{section}}
\newtheorem{example}{Example}[section]
\newenvironment{definitionbox}[1]{%
  \par\noindent\textbf{Definition (#1).}\ \itshape
}{\par\vspace{0.3em}}
\title{File Formats and Named Technology Associations --- Detailed Lecture Notes}
\author{JKSSB Computer Awareness}
\date{\today}

\begin{document}
\maketitle

\section{Extension and format are two different things}
Every file stored on a computer follows some fixed plan for arranging its
data. That plan is the \textbf{format} of the file. The format decides which
program can open the file and what the file is able to store: a picture, a
document, a table, or a database. An \textbf{image format} stores a picture;
a \textbf{document format} stores a printable page; a \textbf{plain-text
format} stores only characters.

A file name usually ends with an \textbf{extension} --- the letters written
after the last dot. In the name \texttt{report.png}, the extension is
\texttt{png}. The extension is only a label; it tells the reader which format
is expected inside. The extension and the format are therefore related but
not the same thing. The exam tests this distinction directly, so it is stated
here in its own section.

\begin{definitionbox}{File format}
A fixed way of arranging the data in a file so that a particular program can
read it, write it, and display it correctly.
\end{definitionbox}

\begin{definitionbox}{File extension}
The letters after the last dot in a file name, such as \texttt{png} in
\texttt{report.png}. The extension signals the file's format.
\end{definitionbox}

\begin{center}
\begin{tabular}{p{0.26\textwidth}p{0.62\textwidth}}
\toprule
\textbf{Term} & \textbf{Meaning in this lesson} \\
\midrule
Extension & The letters after the last dot in a file name. \\
Format & The actual arrangement of data inside the file. \\
PNG & Portable Network Graphics, an image format. \\
PDF & Portable Document Format, a fixed-layout document format. \\
CSV & Comma-Separated Values, a table stored as plain text. \\
TXT & A plain text file with no formatting. \\
\texttt{.mdb} / \texttt{.accdb} & Older and newer Microsoft Access database formats. \\
\bottomrule
\end{tabular}
\end{center}

\section{Image formats}
The best-known image format for this course is \textbf{PNG}, which stands for
\textbf{Portable Network Graphics}. PNG is a raster image format: it stores a
picture as a grid of pixels. It uses \emph{lossless} compression, meaning the
picture is saved without permanently throwing away detail, and it supports
transparency.

Other common image formats are \textbf{JPEG} (written \texttt{.jpg} or
\texttt{.jpeg}), a \emph{lossy} format often used for photographs, and
\textbf{GIF} (\texttt{.gif}), which also supports simple animation. Typical
image extensions are \texttt{.png}, \texttt{.jpg}, \texttt{.jpeg},
\texttt{.gif}, \texttt{.bmp}, and \texttt{.tiff}. The key exam point is
simple: PNG, JPEG, and GIF are all \textbf{image formats}, and none of them
is a document, spreadsheet, or database format.

\begin{definitionbox}{Raster image}
An image stored as a grid of pixels. PNG and JPEG are raster formats.
\end{definitionbox}

\section{PDF: the Portable Document Format}
\textbf{PDF} stands for \textbf{Portable Document Format}, a format created
by Adobe. Its defining property is a \textbf{fixed layout}: a PDF keeps the
exact page layout, fonts, and formatting on every device, so the document
looks the same whether it is opened on a desktop, a phone, or a printer. This
is why PDF is the standard format for sharing a finished document that must
not shift or reflow.

A PDF opens in a \textbf{PDF reader}, the best-known being \textbf{Adobe
Acrobat Reader}. A \textbf{web browser} can also display a PDF, but reading a
PDF is not the same as editing an editable word-processing document. A PDF
may contain real, selectable text, or it may be a \textbf{scanned image} of
each page. A scanned PDF is only a picture; to make its text searchable and
selectable, \textbf{OCR} (Optical Character Recognition) must be applied.
\textbf{PDF/A} is an archival PDF standard meant for long-term preservation.

\begin{example}
A notice is scanned on a photocopier and saved as a PDF. The PDF now contains
pictures of the pages, not typed text. Selecting a sentence with the mouse
does nothing, and searching for a word finds nothing. Running OCR on the file
creates a text layer, after which the same words can be searched and copied.
The format is still PDF; only the internal content has changed from an image
to searchable text.
\end{example}

\section{The Office document family}
Microsoft Office files use a matching family of extensions that are easy to
remember as a group.

\begin{center}
\begin{tabular}{p{0.16\textwidth}p{0.34\textwidth}p{0.36\textwidth}}
\toprule
\textbf{Extension} & \textbf{Format} & \textbf{Program} \\
\midrule
\texttt{.docx} & Word document & Microsoft Word \\
\texttt{.xlsx} & Excel workbook & Microsoft Excel \\
\texttt{.pptx} & PowerPoint presentation & Microsoft PowerPoint \\
\bottomrule
\end{tabular}
\end{center}

These are the modern \textbf{Office Open XML} formats. The older, pre-2007
extensions were \texttt{.doc} (Word), \texttt{.xls} (Excel), and \texttt{.ppt}
(PowerPoint). For an exam, the pairing to hold is: \textbf{.docx = Word},
\textbf{.xlsx = Excel}, \textbf{.pptx = PowerPoint}; each extension always
belongs to the same program, in the modern and the older naming alike.

\section{Plain-data formats: CSV and TXT}
\textbf{CSV} stands for \textbf{Comma-Separated Values}. A CSV file stores a
table as plain text, with the values in each row separated by commas. Because
it is only text, a CSV file can be opened by a spreadsheet program such as
Excel, by a database, or by a simple text editor. The commas are the field
separators; the file itself carries no formatting such as bold text or
colours.

\textbf{TXT} denotes a \textbf{plain text} file. It stores only characters,
with no formatting and no table structure. Both CSV and TXT are simple,
vendor-neutral formats: unlike a \texttt{.docx} or \texttt{.xlsx} file, they
do not belong to one program. The distinction to keep is that \textbf{CSV
holds a table written as text separated by commas}, while \textbf{TXT holds
plain characters with no inherent structure}.

\section{Microsoft Access databases: \texttt{.mdb} and \texttt{.accdb}}
Microsoft Access is the database program of the Office family, and it uses
two extensions that candidates often confuse. \texttt{.mdb} is the
\textbf{older} Access database format, used by Access 2003 and earlier.
\texttt{.accdb} is the \textbf{newer} Access database format, used from
Access 2007 onward. Both files can hold the standard database objects:
tables, queries, forms, reports, macros, and modules.

The important point is that \textbf{both} \texttt{.mdb} and \texttt{.accdb}
are \textbf{Microsoft Access database} files. Neither is a Word document, an
Excel workbook, or a PowerPoint presentation.

\section{Named technology associations}
This section is the high-yield part of the lesson: matching a named
technology to the organisation that produced it or the purpose it serves.

\textbf{Firefox} is a \textbf{web browser} developed by \textbf{Mozilla}.
Mozilla is the organisation, and Firefox is its browser product; Mozilla also
produces \textbf{Thunderbird}, an email client. Firefox must not be confused
with Microsoft's browser, \textbf{Edge} (or the older Internet Explorer), nor
with a PDF reader.

\textbf{C-DAC} stands for the \textbf{Centre for Development of Advanced
Computing}. It is an Indian research and development organisation under the
\textbf{Ministry of Electronics and Information Technology (MeitY)}.
\textbf{PARAM} is a family of \textbf{supercomputers} developed by C-DAC;
\textbf{PARAM 8000} was India's first indigenous supercomputer. The pairing
to remember is \textbf{PARAM = Indian supercomputer, C-DAC = the organisation
that developed it}.

\section{Software roles: browser, PDF reader, and remote support}
Software is often tested by the role it performs rather than by its maker.

\begin{center}
\begin{tabular}{p{0.30\textwidth}p{0.56\textwidth}}
\toprule
\textbf{Role} & \textbf{Example} \\
\midrule
Web browser & Firefox, Google Chrome, Microsoft Edge \\
PDF reader & Adobe Acrobat Reader \\
Remote-support tool & TeamViewer, AnyDesk, Chrome Remote Desktop \\
\bottomrule
\end{tabular}
\end{center}

A \textbf{web browser} is used to visit web pages; a \textbf{PDF reader} is
used to open PDF documents; a \textbf{remote-support tool} lets one computer
control or view another over a network. Firefox is a browser and Adobe
Acrobat Reader is a PDF reader; neither is a remote-support tool. Match the
software to its \emph{job}, not to a similar-sounding name.

\begin{example}
A help-desk officer needs to see a distant user's screen to fix a problem.
The right tool is a remote-support program such as TeamViewer or AnyDesk,
\emph{not} a browser and \emph{not} a PDF reader. If the user then sends a
manual as a fixed-layout file, it will arrive as a PDF and open in a PDF
reader. If the user sends a spreadsheet of data, it will be an XLSX (Excel)
or a CSV. Each task has its own format and its own tool.
\end{example}

\section{Exam retrieval and common confusions}
Six distinctions prevent most errors on this topic.
\begin{enumerate}
  \item An \textbf{extension} is the letters after the dot; a \textbf{format}
    is how the data is stored. ``Extension = format'' is wrong (TRM-09,
    TRAP-34).
  \item \textbf{PNG, JPEG, and GIF are image formats}. CSV is \textbf{not} an
    image format; it is a plain-text table format.
  \item \textbf{PDF} is a fixed-layout document format. A scanned PDF is a
    picture and needs \textbf{OCR} to become searchable.
  \item \textbf{.docx = Word}, \textbf{.xlsx = Excel}, \textbf{.pptx =
    PowerPoint}; the extension always matches the program.
  \item \textbf{.mdb} and \textbf{.accdb} are \textbf{both Access database}
    files --- older and newer respectively.
  \item \textbf{Firefox is a browser by Mozilla}; \textbf{PARAM is a
    supercomputer by C-DAC}. Do not swap the makers.
\end{enumerate}

\subsection*{Exercises}
\begin{enumerate}[label=\arabic*.]
  \item Distinguish a file \emph{extension} from a file \emph{format}, using
    \texttt{report.png} as the example.
  \item Name any three image formats and state one property of PNG.
  \item Explain what is meant by a fixed-layout format, and name the format
    created by Adobe that behaves this way.
  \item Match each extension to its program: \texttt{.docx},
    \texttt{.xlsx}, \texttt{.pptx}, \texttt{.accdb}.
  \item What does CSV stand for, and how is its data arranged?
  \item (a) Who develops Firefox? (b) What is PARAM, and which organisation
    develops it?
\end{enumerate}

\subsection*{Solutions}
\begingroup\small
\begin{enumerate}[label=\arabic*.]
  \item The extension is the suffix after the last dot --- here
    \texttt{png}. The format is the way the data inside the file is actually
    stored --- here PNG, the Portable Network Graphics image format. The
    extension signals the format but is not the format itself.
  \item PNG (Portable Network Graphics), JPEG (\texttt{.jpg/.jpeg}), and GIF
    (\texttt{.gif}), among others. PNG uses lossless compression and supports
    transparency.
  \item A fixed-layout format keeps the exact page layout, fonts, and
    formatting on every device, so the document looks the same wherever it is
    opened. PDF (Portable Document Format) is the format created by Adobe
    that behaves this way.
  \item \texttt{.docx} $\rightarrow$ Microsoft Word; \texttt{.xlsx}
    $\rightarrow$ Microsoft Excel; \texttt{.pptx} $\rightarrow$ Microsoft
    PowerPoint; \texttt{.accdb} $\rightarrow$ Microsoft Access (newer
    database format).
  \item CSV stands for Comma-Separated Values. Its data is stored as a table
    in plain text, with the values in each row separated by commas.
  \item (a) Firefox is developed by Mozilla. (b) PARAM is a family of
    supercomputers, developed by C-DAC (the Centre for Development of
    Advanced Computing); PARAM 8000 was India's first indigenous
    supercomputer.
\end{enumerate}
\endgroup
\end{document}
